% supplementary_material.tex
%
% Build with:  make paper       (from the repository root)
% or:          python3 paper/make_figures.py && pdflatex supplementary_material.tex (x2)
%
% Every figure in this document is regenerated from the pyCICY source by
% paper/make_figures.py. The numbers quoted in the text are written to
% figures/facts.json by that same script, so the prose and the plots cannot
% drift apart silently.

\documentclass[11pt,a4paper]{article}

\usepackage[T1]{fontenc}
\usepackage[utf8]{inputenc}
\usepackage[margin=1in]{geometry}
\usepackage[protrusion=true,expansion=false]{microtype}
\usepackage{amsmath,amssymb,amsthm}
\usepackage{graphicx}
\usepackage{booktabs}
\usepackage{enumitem}
\usepackage{xcolor}
\usepackage[hidelinks]{hyperref}
\usepackage[nameinlink,capitalize]{cleveref}
\usepackage{framed}
\usepackage{hyperref}
\usepackage{orcidlink}


\graphicspath{{figures/}}

% Numbers computed by paper/make_figures.py. Defined here as a fallback so
% the document still typesets if the figures have not been generated yet;
% the \input below overrides them with the real values.
\newcommand{\FactDepth}{?}\newcommand{\FactConfigs}{?}
\newcommand{\FactEdges}{?}\newcommand{\FactEffective}{?}
\newcommand{\FactIneffective}{?}\newcommand{\FactHodgePairs}{?}
\newcommand{\FactFavourable}{?}\newcommand{\FactQuinticNodes}{?}
\newcommand{\FactValidationPairs}{?}\newcommand{\FactValidationAgree}{?}
\newcommand{\FactHminOne}{?}\newcommand{\FactHmaxOne}{?}
\newcommand{\FactHminTwo}{?}\newcommand{\FactHmaxTwo}{?}
\newcommand{\FactCold}{?}\newcommand{\FactWarm}{?}
\newcommand{\FactChTwoConfigs}{?}\newcommand{\FactChTwoCoeffs}{?}
\newcommand{\FactChTwoDisagree}{?}\newcommand{\FactChTwoEdges}{?}
\newcommand{\FactChTwoIneffective}{?}\newcommand{\FactChTwoIneffNonzero}{?}
\newcommand{\FactGVModels}{?}\newcommand{\FactGVMaxDegree}{?}
\newcommand{\FactGVQuinticOne}{?}\newcommand{\FactGVQuinticTwo}{?}
\newcommand{\FactGVQuinticThree}{?}
\newcommand{\FactGenCicys}{?}\newcommand{\FactGenCandidates}{?}
\newcommand{\FactGenWithSymmetry}{?}\newcommand{\FactGenRecords}{?}
\newcommand{\FactGenHits}{?}\newcommand{\FactGenMinEuler}{?}
\newcommand{\FactGenList}{?}
\newcommand{\FactAddChecked}{?}\newcommand{\FactAddMultiple}{?}
\newcommand{\FactAddMismatch}{?}\newcommand{\FactAddSpread}{?}
\makeatletter
\@ifundefined{input@path}{}{}
\makeatother
\IfFileExists{figures/facts.tex}{\input{figures/facts.tex}}{}

\newcommand{\PP}{\mathbb{P}}
\newcommand{\cO}{\mathcal{O}}
\newcommand{\ch}{\operatorname{ch}}
\newcommand{\XD}{X_D}
\newcommand{\XR}{X_R}

\theoremstyle{definition}
\newtheorem{remark}{Remark}

\title{\vspace{-1cm}\textbf{What a Configuration Matrix Determines,\\
and What It Does Not:\\[2pt]
\large Exact Methods for Heterotic Compactifications on CICYs}}
\author{B.S. Hartshorn \orcidlink{0009-0004-2853-655X} \small \url{https://github.com/brentharts/CICY}}

\begin{document}
\maketitle

\begin{abstract}
The organizing question is which quantities in a heterotic compactification
are determined by the configuration matrix alone. A surprising number are, and
they are exact --- integers or integer series, computed symbolically and
cross-checked by independent routes. The node count of a conifold transition,
the second Chern character across it, the index of a line bundle sum, the
character-valued index of a discrete group action, the branching of an
$SU(5)$ Wilson line, the Yukawa coupling of the standard embedding together
with its worldsheet instanton corrections, and the entire pattern of allowed
and forbidden couplings of a line bundle model: all of these follow from the
degrees. An equal number are not, and the boundary is sharp rather than gradual. The
holomorphic Yukawa couplings of a line bundle model are quasi-topological but
need explicit cohomology representatives; the physical couplings need the
Ricci-flat metric; the masses need moduli stabilisation as well. For a freely acting $\Gamma$ the
equivariant index is a multiple of the regular representation, so the choice
of equivariant structure --- $n^r$ of them --- does not affect the chiral
spectrum at all. The ambiguity is confined to the vector-like sector, which an
index cannot see in any case.
\end{abstract}

\medskip

% =====================================================================
\section{Configuration matrices, normal forms and equivalence}
\label{sec:cipro}

A \emph{complete intersection Calabi--Yau manifold} (CICY) is specified by a configuration matrix
%
\begin{equation}
X \;=\;
\begin{bmatrix}
\PP^{n_1} & q^1_1 & \cdots & q^1_K\\
\vdots    & \vdots&        & \vdots\\
\PP^{n_m} & q^m_1 & \cdots & q^m_K
\end{bmatrix},
\qquad q^r_a \in \mathbb{Z}_{\ge 0},
\label{eq:conf}
\end{equation}
%
the zero locus of $K$ multi-homogeneous polynomials in
$\mathcal{A}=\PP^{n_1}\times\cdots\times\PP^{n_m}$. The threefold condition
and the vanishing of the first Chern class read
%
\begin{equation}
\sum_{r=1}^{m} n_r - K = 3,
\qquad
\sum_{a=1}^{K} q^r_a = n_r + 1 \quad \text{for each } r .
\label{eq:cy}
\end{equation}

Different matrices can describe the same variety. The simplest redundancy is
a relabelling: permuting the projective factors permutes the rows of
\cref{eq:conf}, and permuting the defining polynomials permutes its columns.
\texttt{pyCICY} canonicalises under this action, mirroring the
\texttt{CINormal} and \texttt{CIEquiv} functionality of
\texttt{CIPro}~\cite{Anderson2026}. Naively the orbit has $m!\,K!$ elements,
which is impractical at the sizes occurring in the list (up to $12\times 15$).
Instead, rows and columns are partitioned by iterated signature refinement:
each row is labelled by its dimension together with the multiset of
(degree, column label) pairs it contains, and dually for columns, until the
labelling is stable. Only permutations \emph{within} a class can be
symmetries, so the search reduces to a product of factorials of class sizes,
which is small in practice. An explicit budget is enforced rather than
allowing an unbounded run.

\begin{remark}
The canonical representative chosen by \texttt{pyCICY} is its own. It
identifies the same equivalences as \texttt{CIPro} but does not necessarily
select the same matrix from an orbit; only the induced equivalence relation
is claimed to agree.
\end{remark}

% =====================================================================
\section{Splits and conifold transitions}
\label{sec:transitions}

\subsection{The split construction}

A $\PP^1$ split replaces one defining polynomial by two whose degrees sum to
the original, at the cost of one extra $\PP^1$ factor~\cite{Candelas1988}.
Writing the split column as $q$ and choosing a partition
$q = a + b$ row by row,
%
\begin{equation}
\begin{bmatrix} \PP^{n_r} & q^r & \cdots \end{bmatrix}
\;\longmapsto\;
\begin{bmatrix}
\PP^{1}   & 1   & 1   & 0\cdots\\
\PP^{n_r} & a^r & b^r & \cdots
\end{bmatrix} .
\label{eq:split}
\end{equation}
%
Both conditions in \cref{eq:cy} are preserved: the row sums are unchanged
because $a^r + b^r = q^r$, the new row sums to $2 = n+1$, and the ambient
dimension and the number of equations each increase by one. For the quintic
this produces
%
\begin{equation}
\XD=\begin{bmatrix}\PP^4 & 5\end{bmatrix}
\;\longleftrightarrow\;
\begin{bmatrix}\PP^1 & 1 & 1\\ \PP^4 & 1 & 4\end{bmatrix}=\XR ,
\label{eq:quinticsplit}
\end{equation}
%
which is equation (1.11) of~\cite{Anderson2025}. The package reproduces
\cref{eq:quinticsplit} exactly, and contracting $\XR$ returns
$\begin{bmatrix}\PP^4 & 5\end{bmatrix}$.

\subsection{Topology across the transition}

Geometrically the two sides of \cref{eq:split} are joined through a nodal
variety $X_0$: $\XD$ is a deformation of $X_0$ and $\XR$ a small resolution.
Following Anderson and Gray (2025) ~\cite{Anderson2025},
%
\begin{align}
h^{1,1}(\XR) &= h^{1,1}(\XD) + 1, \label{eq:h11}\\
\chi(\XR)    &= \chi(\XD) + 2N,   \label{eq:chi}
\end{align}
%
with $N$ the number of nodes of $X_0$. \Cref{eq:chi} is used to
\emph{measure} $N$, since $\chi$ follows from the triple intersection numbers
and Chern classes that \texttt{pyCICY} already computes, while \cref{eq:h11}
is reported as a check rather than assumed.

\subsection{An independent route to \texorpdfstring{$N$}{N}}
\label{sec:nodesexpected}

There is a second way to obtain $N$ that uses no cohomology whatever. In the
$\PP^1$-split presentation the two defining equations of $\XR$ are
$x_0 f_1 + x_1 f_2$ and $x_0 g_1 + x_1 g_2$, and the nodes of $X_0$ are the
points where all four of $f_1,f_2,g_1,g_2$ vanish together with the remaining
defining equations. Those four polynomials have multidegrees $a,a,b,b$, so
%
\begin{equation}
N \;=\;
\int_{\mathcal{A}}
\;(a\!\cdot\!J)^2 \,(b\!\cdot\!J)^2
\prod_{c \neq \text{split}} (q_c\!\cdot\!J),
\qquad a\!\cdot\!J \equiv \sum_r a^r J_r ,
\label{eq:nodesint}
\end{equation}
%
the intersection number in the ambient space, evaluated by extracting the
coefficient of $\prod_r J_r^{n_r}$ subject to $J_r^{n_r+1}=0$. The count of
classes in \cref{eq:nodesint} is $4 + (K-1) = K+3$, which equals
$\dim\mathcal{A}=\sum_r n_r$ precisely when \cref{eq:cy} holds, so the result
is a number and not a class. For the split of \cref{eq:quinticsplit},
$a=1$ and $b=4$ in $\PP^4$ give $1\cdot1\cdot4\cdot4=16$, and the partition
$(2,3)$ gives $2\cdot2\cdot3\cdot3=36$.

\cref{eq:nodesint} is a coefficient
extraction over a truncated polynomial ring, whereas \cref{eq:chi} runs
through Bott--Borel--Weil, the Koszul resolution and the Leray spectral
sequence. An error in either would show up as a point off the diagonal.

For the transition of \cref{eq:quinticsplit} both routes give $N=\FactQuinticNodes$, Anderson 2025 ~\cite{Anderson2025}.

\subsection{The second Chern character across the transition}
\label{sec:chtwo}

The relation quoted as \cref{eq:ch2} in \cref{sec:aside} can be checked
directly. From the Euler sequence for the ambient space and the adjunction
sequence for the complete intersection,
%
\begin{equation}
\ch(T\mathcal{A}) = \sum_r \Bigl[(n_r+1)\,e^{J_r} - 1\Bigr],
\qquad
\ch(TX) = \ch(T\mathcal{A})\big|_X - \sum_a \ch\bigl(\cO(q_a)\bigr),
\label{eq:chadj}
\end{equation}
%
whose degree-two part is
%
\begin{equation}
\ch_2(TX) \;=\; \tfrac{1}{2}\Bigl[\;
\sum_r (n_r+1)\,J_r^{\,2}
\;-\;\sum_a (q_a\!\cdot\!J)^2 \;\Bigr],
\qquad J_r^{\,n_r+1} = 0 .
\label{eq:chtwo}
\end{equation}

To compare the two sides of a transition their classes must live in one
ambient space. Following Anderson and Gray (2025) eq.~(1.29) of~\cite{Anderson2025}, the deformation
is re-embedded in the ambient space of the resolution by cutting the new
$\PP^1$ down to a point,
%
\begin{equation}
\begin{bmatrix}\PP^4 & 5\end{bmatrix}
\;\longmapsto\;
\begin{bmatrix}\PP^1 & 1 & 0\\ \PP^4 & 0 & 5\end{bmatrix},
\label{eq:reembed}
\end{equation}
%
which describes $\{\text{pt}\}\times X_D \cong X_D$. \Cref{eq:chtwo} then
gives, for the split of \cref{eq:quinticsplit},
%
\begin{equation}
\ch_2(T\XD) = -10 D_2^2,
\qquad
\ch_2(T\XR) = -5D_1D_2 - 6D_2^2,
\qquad
\bigl[\PP^1_s\bigr] = -5D_1D_2 + 4D_2^2 ,
\label{eq:ch124}
\end{equation}
%
reproducing eq.~(1.24) of~\cite{Anderson2025} term by term. The bridging
curves of \cref{eq:bridging} come out as $[C_D] = 4D_2^2$ and
$[C_R] = 5D_1D_2$, also as quoted there.

\begin{remark}[What is and is not being tested]
\label{rem:identity}
Writing $q = a + b$ and $[C_D] = (a\!\cdot\!J)(b\!\cdot\!J)$,
$[C_R] = J_0\,(q\!\cdot\!J)$, one finds
%
\begin{equation}
\ch_2(T\XR) - \ch_2(T\XD)
= (a\!\cdot\!J)(b\!\cdot\!J) - J_0\,(q\!\cdot\!J)
= [C_D] - [C_R]
\label{eq:identity}
\end{equation}
%
term by term, directly from \cref{eq:chtwo}. \Cref{eq:identity} is therefore
an \emph{algebraic identity}, not an independent check: it cannot fail once
\cref{eq:chtwo} and the re-embedding are adopted. It is still evaluated, as a
guard against an implementation error in the class arithmetic, but it is not
evidence that \cref{eq:chtwo} is right. The substantive check is different:
for a Calabi--Yau, $c_1=0$ gives $\ch_2 = -c_2$, and \texttt{pyCICY}
computes $c_2^{rs}$ by an unrelated route~\cite{Larfors2019}, ~\cite{SuppFigures}.
\end{remark}

The companion supplement~\cite{SuppFigures} reports both (Figure~1). The right-hand panel records a point that is
easy to get wrong: a nonzero exceptional class does not imply that the
transition does anything. For an ineffective split the nodal locus is empty,
$N=0$ by \cref{eq:nodesint}, and yet $[\PP^1_s]$ computed from
\cref{eq:identity} is a perfectly good nonzero class. The two quantities
answer different questions, and only the intersection number answers
``how many nodes''.

% =====================================================================
\section{The split web}
\label{sec:web}

\subsection{Construction}

The classification of~\cite{Candelas1988} was built by splitting: starting
from the configurations with a single defining equation and applying
\cref{eq:split} repeatedly, discarding duplicates~\cite{Candelas1988,Green1987, Green1989}. For $K=1$ the threefold
condition in \cref{eq:cy} forces $\sum_r n_r = 4$ and then the Calabi--Yau
condition fixes the single column to $q^r = n_r+1$, so there are exactly five
seeds:
%
\begin{equation}
\begin{bmatrix}\PP^4 & 5\end{bmatrix},\quad
\begin{bmatrix}\PP^3 & 4\\ \PP^1 & 2\end{bmatrix},\quad
\begin{bmatrix}\PP^2 & 3\\ \PP^2 & 3\end{bmatrix},\quad
\begin{bmatrix}\PP^2 & 3\\ \PP^1 & 2\\ \PP^1 & 2\end{bmatrix},\quad
\begin{bmatrix}\PP^1 & 2\\ \PP^1 & 2\\ \PP^1 & 2\\ \PP^1 & 2\end{bmatrix}.
\label{eq:seeds}
\end{equation}
%
Breadth-first splitting from \cref{eq:seeds}, reducing every configuration to
the normal form of \cref{sec:cipro} before recording it, reaches $\FactConfigs$
configurations at depth $\FactDepth$, joined by $\FactEdges$ split edges, of
which $\FactEffective$ are effective and $\FactIneffective$ have $N=0$.

\subsection{Ineffective splits}

Splits with $N=0$ are the \emph{ineffective} splits of~\cite{Candelas1988}:
the would-be nodal locus is empty, no topology changes, and the two
configurations describe isomorphic geometries. The example of section~5.2
of~\cite{Anderson2025},
%
\begin{equation}
\begin{bmatrix}\PP^1 & 2\\ \PP^3 & 4\end{bmatrix}
\;\longleftrightarrow\;
\begin{bmatrix}\PP^1 & 1 & 1\\ \PP^1 & 1 & 1\\ \PP^3 & 4 & 0\end{bmatrix},
\label{eq:ineffective}
\end{equation}
%
is reproduced: $h^{1,1}$, $h^{2,1}$ and $\chi$ are all unchanged, and
\cref{eq:nodesint} returns $N=0$. Note that \cref{eq:h11} must therefore
\emph{not} be imposed unconditionally; the expected shift in $h^{1,1}$ is $1$
for an effective transition and $0$ for an ineffective one. The two
descriptions in \cref{eq:ineffective} differ in favourability, the right-hand
side being unfavourable, which is precisely the mechanism by which
ineffective splitting is used in the literature to move between favourable
and unfavourable presentations of one manifold.

\subsection{The survey}

%% figure moved to the companion supplement \cite{SuppFigures}


The companion supplement~\cite{SuppFigures} presents the survey (Figures 1-A and 3-A).
All configurations generated satisfy \cref{eq:cy}, have no degenerate rows or
columns, and have $\chi \le 0$, as every CICY threefold must; the same web
underlies~\cite{AndersonGaoGrayLee2017,AndersonFibrations2018}. The observed
ranges $h^{1,1}\in[\FactHminOne,\FactHmaxOne]$ and
$h^{2,1}\in[\FactHminTwo,\FactHmaxTwo]$ sit inside the published $[0,19]$ and
$[0,101]$.

\Cref{tab:landmarks} lists manifolds that can be checked against the
literature. The Schoen manifold is worth singling out: it is not something
the construction was aimed at, but arises as the split bicubic
$\left[\begin{smallmatrix}\PP^1 & 1 & 1\\ \PP^2 & 0 & 3\\
\PP^2 & 3 & 0\end{smallmatrix}\right]$, and comes out with
$h^{1,1}=h^{2,1}=19$ and $\chi=0$, its known values.

\begin{table}[t]
\centering
\begin{tabular}{llrrr}
\toprule
Manifold & Configuration & $h^{1,1}$ & $h^{2,1}$ & $\chi$ \\
\midrule
Quintic          & $[\PP^4\,|\,5]$                       & 1  & 101 & $-200$\\
Split quintic    & \cref{eq:quinticsplit}                & 2  & 86  & $-168$\\
Bicubic          & $[\PP^2\,|\,3;\ \PP^2\,|\,3]$          & 2  & 83  & $-162$\\
Tetraquadric     & $[\PP^1\,|\,2]^{\times 4}$             & 4  & 68  & $-128$\\
Schoen           & split bicubic                          & 19 & 19  & $0$\\
\bottomrule
\end{tabular}
\caption{Landmark manifolds appearing in the generated web, with the values
\texttt{pyCICY} computes. All five agree with the literature. The first four
are used as regression tests; the fifth was not targeted and emerged from the
splitting construction of \cref{sec:web}.}
\label{tab:landmarks}
\end{table}

% =====================================================================
\section{Chern polynomials, Hilbert series and enumerative invariants}
\label{sec:enumerative}

The quantities of this section are self-contained, in that each can be
checked against a published value without reference to anything else in this
document. They correspond to section~1 of~\cite{Anderson2026}.

\subsection{Chern polynomial}

The tangent sequence $0 \to TX \to T\mathcal{A}|_X \to \mathcal{E}|_X \to 0$
with $\mathcal{E} = \bigoplus_a \cO_{\mathcal{A}}(q_a)$ gives
%
\begin{equation}
c(TX) \;=\; \frac{c(T\mathcal{A})}{c(\mathcal{E})}
\;=\; \frac{\prod_i (1+J_i)^{n_i+1}}{\prod_a (1 + q_a\!\cdot\!J)} ,
\label{eq:cpoly}
\end{equation}
%
truncated by $J_i^{n_i+1}=0$. For the $dP_2$ surface of~\cite{Anderson2026}
this returns $1 + J_1 + J_2 + 2J_1J_2 + J_3 + 2J_1J_3 + J_2J_3$, as
published. For the quintic it gives $c(TX) = 1 + 10J^2 - 40J^3$, so that
%
\begin{equation}
\chi \;=\; \int_X c_3(TX) \;=\; \int_{\mathcal{A}} c_3(TX)\,\prod_a
(q_a\!\cdot\!J) \;=\; -40 \times 5 \;=\; -200 .
\label{eq:eulerc3}
\end{equation}
%
\Cref{eq:eulerc3} is a route to $\chi$ independent of the triple intersection
numbers \texttt{pyCICY} normally uses, and the two agree on every
configuration of the web. The degree-two part of \cref{eq:cpoly} likewise
reproduces $-\ch_2$ from \cref{eq:chtwo}, which ties this section back to
\cref{sec:chtwo}.

\subsection{Hilbert series}

The Koszul resolution of $\cO_X$ by the defining equations gives the
multigraded Hilbert series of the coordinate ring as
%
\begin{equation}
H(t) \;=\;
\frac{\prod_a \bigl(1 - t^{q_a}\bigr)}{\prod_i (1-t_i)^{n_i+1}},
\qquad t^{q_a} = \prod_i t_i^{\,q^i_a} .
\label{eq:hilbert}
\end{equation}
%
For CICY $7821$, the configuration
$\left[\begin{smallmatrix}\PP^2 & 1&1&1\\ \PP^4 & 2&2&1\end{smallmatrix}
\right]$ of~\cite{Anderson2026}, \cref{eq:hilbert} gives
%
\begin{equation}
H(t) \;=\; \frac{(1-t_1t_2)(1-t_1t_2^2)^2}{(1-t_1)^3(1-t_2)^5},
\label{eq:hs7821}
\end{equation}
%
character for character the published expression. The same manifold has
$h^{1,1}=2$, $h^{2,1}=58$ and $\chi=-112$ in~\cite{Anderson2026}, and
\texttt{pyCICY} returns those values, so the configuration is anchored on
both sides.

\subsection{Gopakumar--Vafa invariants}
\label{sec:gv}
\label{sec:gv}

Genus-zero GV invariants are implemented for the \emph{one-parameter} models
only: the $\FactGVModels$ Calabi--Yau threefolds that are complete
intersections in a single projective space,
%
\begin{equation}
\PP^4[5],\quad \PP^5[3,3],\quad \PP^5[2,4],\quad \PP^6[2,2,3],\quad
\PP^7[2,2,2,2] .
\label{eq:oneparam}
\end{equation}
%
These are exactly the CICYs with a single K\"ahler parameter, which is what
makes the mirror map a one-variable problem. With degrees $d_a$ in $\PP^N$,
write $\kappa = \prod_a d_a$ and $\mu = \prod_a d_a^{d_a}$. The fundamental
period is
%
\begin{equation}
w_0(z) \;=\; \sum_{k\ge 0}
\frac{\prod_a (d_a k)!}{(k!)^{N+1}}\, z^k ,
\label{eq:period}
\end{equation}
%
and the logarithmic solution supplies the mirror map $t = \log z +
\sigma/w_0$, whose coefficients involve
$\sum_a d_a H_{d_a k} - (N+1) H_k$ with $H$ the harmonic numbers. The
Euler--Mascheroni constants cancel from that bracket precisely because
$\sum_a d_a = N+1$, which is the Calabi--Yau condition \cref{eq:cy}. The
Yukawa coupling
%
\begin{equation}
K_{ttt} \;=\; \kappa \, + \sum_{d \ge 1} n_d\, \frac{d^3 q^d}{1-q^d},
\qquad q = e^{t},
\label{eq:yukawa}
\end{equation}
%
then yields the invariants $n_d$. The method is that
of Hosono, et al. (1994) ~\cite{HKTY1994}; \texttt{CIPro} obtains its GV invariants from
code by A.~Klemm based on the same references, and what is described here is
an independent implementation restricted to \cref{eq:oneparam}, not a port.

All five reproduce the tabulated instanton numbers to degree
$\FactGVMaxDegree$, the quintic beginning $\FactGVQuinticOne$,
$\FactGVQuinticTwo$, $\FactGVQuinticThree$; integrality is itself the check, since \cref{eq:yukawa}
has no reason to yield integers unless the mirror map and Yukawa coupling are
both right, and it is enforced rather than assumed. Multi-parameter
models require the full multivariable mirror map and are refused with an
explicit error rather than silently mishandled; see \cref{sec:limits}.

% =====================================================================
\section{From topology to generations}
\label{sec:generations}

\subsection{What the configuration matrix determines}

For the heterotic string with the standard embedding $V = TX$ the gauge group
is $E_6$ and the charged spectrum is counted by Hodge numbers: $n(\mathbf{27})
= h^{2,1}$, $n(\overline{\mathbf{27}}) = h^{1,1}$, so the net number of chiral
generations is
%
\begin{equation}
n_{\rm gen} \;=\; h^{2,1} - h^{1,1} \;=\; -\tfrac{1}{2}\chi(X) ,
\label{eq:gen}
\end{equation}
%
the index of the Dirac operator twisted by $V$. On a quotient by a freely
acting $\Gamma$ this divides,
%
\begin{equation}
n_{\rm gen}\bigl(X/\Gamma\bigr) \;=\; \frac{|\chi(X)|}{2|\Gamma|} ,
\label{eq:genquot}
\end{equation}
%
so three generations requires $|\chi| = 6|\Gamma|$.

Across the published list the smallest non-zero $|\chi|$ is
$\FactGenMinEuler$, and $|\chi| = 6$ never occurs. A three-generation
standard-embedding model therefore \emph{always} requires a non-trivial free
quotient --- which is convenient, since the quotient is also what supplies the
Wilson lines needed to break $E_6$ towards the Standard Model group.

\subsection{What it does not determine}
\label{sec:notdetermined}

Masses and couplings are not functions of the configuration matrix, and no
amount of computation on it will produce them. Physical Yukawa couplings are
holomorphic couplings divided by the K\"ahler normalisation of the fields;
that normalisation requires the Ricci-flat metric, for which no closed form is
known on any compact Calabi--Yau threefold. The overall scale requires the
supersymmetry breaking mechanism, the gauge coupling requires the dilaton
vacuum expectation value, and all of these require the moduli to be
stabilised, which is unsolved across the landscape.

\subsection{Narrowing by the actual symmetries}
\label{sec:narrowing}

\Cref{eq:genquot} says which $|\Gamma|$ \emph{would} be needed; it does not
say that a freely acting group of that order exists. That is the content
of~\cite{Braun2010}, whose classification ships with the Mathematica version
of the CICY list and records each symmetry by its \textsc{Gap} SmallGroup
identifier. Hodge numbers for the resulting quotients were computed
in~\cite{ConstantinGrayLukas2017,CandelasConstantinMishra2016}, and the
classification has since been extended to the favourable
list~\cite{GrayWang2022}. Reading it gives $\FactGenWithSymmetry$ configurations carrying
symmetry data and $\FactGenRecords$ records in total, matching the counts
quoted for that classification in~\cite{Braun2010,LukasMishra2020}.

%% figure moved to the companion supplement \cite{SuppFigures}


The companion supplement~\cite{SuppFigures} shows the effect (Figure 3). Of the $\FactGenCandidates$
configurations admitted by \cref{eq:genquot}, exactly $\FactGenHits$ carry a
recorded group of the required order, namely CICYs $\FactGenList$. The first
of these is
%
\begin{equation}
X \;=\;
\begin{bmatrix}\PP^3 & 3 & 0 & 1\\ \PP^3 & 0 & 3 & 1\end{bmatrix},
\qquad h^{1,1} = 14, \quad h^{2,1} = 23, \quad \chi = -18,
\label{eq:tianyau}
\end{equation}
%
with a freely acting $\mathbb{Z}_3$, so that $\chi(X/\mathbb{Z}_3) = -6$ and
$n_{\rm gen} = 3$. \Cref{eq:tianyau} is the Tian--Yau
manifold~\cite{TianYau1987,CandelasIII}, the first three-generation
Calabi--Yau constructed. It was not put into the search but came out of it,
exercising the list parse, the \textsc{Gap} extraction and \cref{eq:gen} together.

\begin{remark}[The record flag, and why it does not change the count]
\label{rem:flag}
Each symmetry record carries a leading boolean which splits them $1199$ to
$496$. Restricting to the records where it is true would leave only the
Tian--Yau manifold of \cref{eq:tianyau} rather than $\FactGenHits$
configurations, so it matters what it means.

It is \emph{not} freeness. The literature is explicit that all
$\FactGenRecords$ symmetries in this classification are freely acting: of the
$7890$ CICY threefolds, $\FactGenWithSymmetry$ exhibit known discrete
symmetries, and those $\FactGenWithSymmetry$ configuration matrices carry
$\FactGenRecords$ symmetries between them, giving $\FactGenRecords$ quotient
manifolds~\cite{Braun2010,ConstantinGrayLukas2017,AndersonFibrations2018,%
LukasMishra2020}. Restricting to the flagged records is therefore not
justified, and the count of three-generation candidates stands at
$\FactGenHits$.

What the flag does mark was not determined here, though the data constrains
it. Every one of the $1199$ true-flagged records has prime-power order, and
every record whose order is not a prime power --- orders $6$, $10$, $12$ and
$20$, $41$ records in all --- is false-flagged. The converse fails: $455$
records of prime-power order are also false-flagged, so the flag is not
simply the condition of being a $p$-group. It is not the divisibility
$|\Gamma| \mid \chi$ either, which holds for all $\FactGenRecords$ records on
both sides, nor is it a property of the configuration, since $40$
configurations carry records of both kinds while $79$ carry only true and
$76$ only false.
\end{remark}

\begin{remark}[Consistency of the source data]
The \textsc{Gap} order can be checked against the abelian invariants recorded
alongside it, and the two agree in $1630$ of the $1631$ abelian cases. The
exception carries order $8$ against invariants $\{2,8\}$, whose product is
$16$. That is an inconsistency in the published data; both readings are
surfaced and neither is silently preferred.
\end{remark}

% =====================================================================
\section{Additivity along composed splits}
\label{sec:additivity}

Brittenham and Hermiller~\cite{Brittenham2025} have shown that unknotting
number is not additive under connected sum, exhibiting
%
\begin{equation}
u\bigl(7_1 \# \overline{7_1}\bigr) \;\le\; 5 \;<\; 6
\;=\; u(7_1) + u(\overline{7_1}),
\label{eq:unknotting}
\end{equation}
%
which settles problem 1.69(B) of Kirby's list negatively \cite{Wang2025}. Unknotting number
is a \emph{minimal move count}: the fewest crossing changes carrying a knot
to the unknot. The content of \cref{eq:unknotting} is that the complexity of
a composite is not determined by that of its factors.

The split web of \cref{sec:web} has the same shape: a family of objects, an
elementary move, and distinguished trivial objects, namely the five
single-equation seeds of \cref{eq:seeds}. It is therefore natural to ask
whether the analogous complexity measures misbehave in the analogous way.
They do not, and the reason is worth stating because it is structural rather
than accidental.

\subsection{What cannot vary}

The direct analogue of unknotting number is the \emph{split depth}: the
fewest splits carrying a seed to a given configuration. It is rigid. A split
adds exactly one defining equation, a contraction removes exactly one, and
the seeds are precisely the configurations with $K=1$, so
%
\begin{equation}
\text{split depth} \;=\; K - 1
\label{eq:depth}
\end{equation}
%
whatever route is taken. Every configuration in the web satisfies
\cref{eq:depth}; every contraction path has that same length; greedy
contraction is automatically optimal; and no contraction path terminates
anywhere but at a seed. There is accordingly no room here for the phenomenon
of \cref{eq:unknotting}, nor for a Bernhard--Jablan style failure of greedy
descent~\cite{Brittenham2021}. The package verifies each of these rather than
inferring them from \cref{eq:depth}.

The total node count is rigid for a different reason. Along any chain of
splits from $\XD$ to $\XR$ the individual counts satisfy
%
\begin{equation}
\sum_i N_i \;=\; \frac{\chi(\XR) - \chi(\XD)}{2},
\label{eq:telescope}
\end{equation}
%
since the Euler characteristics telescope by \cref{eq:chi}. The right-hand
side depends only on the endpoints, so the total cannot depend on the route.

\subsection{What does vary}

The \emph{decomposition} of that total is not fixed. Of the
$\FactAddChecked$ non-seed configurations surveyed, $\FactAddMultiple$ admit
more than one multiset of per-step node counts. For example the configuration
$\left[\begin{smallmatrix}\PP^1&0&1&1\\ \PP^1&1&0&1\\
\PP^4&1&2&2\end{smallmatrix}\right]$ is reached both through
$(N_1,N_2) = (12,36)$ and through $(16,32)$; both sum to $48$, as
\cref{eq:telescope} requires. Individual steps differ by as much as
$\FactAddSpread$ between routes, while the totals disagree in
$\FactAddMismatch$ cases out of $\FactAddChecked$ --- which is also a check on
\cref{eq:chi} across the whole web, since a single arithmetic slip anywhere
in the cohomology would break the telescoping.

The common to both settings is the one worth carrying away: a quantity that looks as though it should be
assembled additively out of elementary steps need not be, and which of the
two failure modes occurs --- an undetermined total, or an undetermined
decomposition --- is a fact about the structure rather than something to be
assumed. In the case at hand the exceptional class of \cref{sec:chtwo} is a
further instance of the same caution, being an invariant that looks as though
it should detect a property it does not.

\begin{remark}
This section takes a question from~\cite{Brittenham2025} and asks it here; it
does not transfer a theorem. The two settings share no formalism, and nothing
in \cref{sec:transitions,sec:web} depends on the comparison. What is claimed
is only what was computed: \cref{eq:depth,eq:telescope} hold across the
generated web, and the decomposition varies.
\end{remark}

% =====================================================================
\section{Bundles beyond the standard embedding}
\label{sec:bundles}

An exact sign obstruction to poly-stability: at a common zero of all the slopes every \emph{partial} sum of
slopes vanishes too, so a subset whose summed slope form has no negative entry
rules the bundle out --- $2^n-2$ sign tests replacing a numerical search that
is a thousand times slower. 

\Cref{sec:generations} treats $V = TX$, where the spectrum is a function of the
Hodge numbers and nothing has to be chosen. That is the exception. A general
heterotic model requires a holomorphic bundle $V$, and the bundle is extra
data: not a function of the configuration matrix. What \emph{is} a function of
the configuration matrix, once $V$ is given, is everything an index theorem can
say about it.

\subsection{Sums of line bundles}

Take
%
\begin{equation}
V \;=\; \cO(L_1) \oplus \cdots \oplus \cO(L_n),
\qquad \sum_{a=1}^{n} L_a \;=\; 0 ,
\label{eq:lbsum}
\end{equation}
%
so that the structure group is $S\bigl(U(1)^n\bigr) \subset SU(n)$. These are
the bundles of~\cite{AndersonGrayLukasPalti2011}, and they are the tractable
case because cohomology commutes with direct sums: $h^i(V) = \sum_a h^i(L_a)$,
each summand computable exactly.

On a threefold with $c_1(V) = 0$ the index theorem gives
$\operatorname{ind}(V) = \int_X \ch_3(V)$, and for \cref{eq:lbsum} this is
%
\begin{equation}
\operatorname{ind}(V)
\;=\; \frac{1}{6}\, d_{rst} \sum_{a} L_a^r L_a^s L_a^t ,
\qquad n_{\rm gen} \;=\; -\operatorname{ind}(V) ,
\label{eq:lbindex}
\end{equation}
%
in the same convention as \cref{eq:gen}: for $V = TX$ one has
$\ch_3 = c_3/2$, so $\operatorname{ind}(TX) = \chi/2$ and $n_{\rm gen} =
-\chi/2$, recovering \cref{eq:gen} exactly.

\Cref{eq:lbindex} admits a second derivation which shares no code with the
first, and the reason it does is worth displaying. The Atiyah--Singer index of
a single line bundle carries a subleading term,
%
\begin{equation}
\operatorname{ind}\bigl(\cO(L)\bigr)
\;=\;
\underbrace{\tfrac{1}{6}\, d_{rst} L^r L^s L^t}_{\text{cubic}}
\;+\;
\underbrace{\tfrac{1}{12}\, c_2^{\,r} L_r}_{\text{linear in } L} ,
\label{eq:lineindex}
\end{equation}
%
and summing \cref{eq:lineindex} over the summands of \cref{eq:lbsum} gives
%
\begin{equation}
\operatorname{ind}(V)
\;=\;
\tfrac{1}{6}\, d_{rst} \sum_a L_a^r L_a^s L_a^t
\;+\;
\tfrac{1}{12}\, c_2^{\,r}
\overbrace{\Bigl(\textstyle\sum_a L_a\Bigr)_{\!r}}^{=\,0
\ \text{by \cref{eq:lbsum}}} ,
\end{equation}
%
so the linear term drops out precisely because the structure group is
$SU(n)$ rather than $U(n)$. The two routes --- one contracting triple
intersection numbers, the other running the Leray spectral sequence on each
summand --- agree identically, and their agreement is a genuine check rather
than a tautology.

\subsection{Anomaly cancellation, and why it is not vacuous}

The Bianchi identity requires $c_2(TX) - c_2(V) - c_2(\tilde V) = [W]$ with
$[W]$ effective. For \cref{eq:lbsum} with $c_1(V) = 0$,
%
\begin{equation}
c_2(V) \;=\; -\ch_2(V) ,
\qquad
\ch_2(V)_r \;=\; \tfrac{1}{2}\, d_{rst} \sum_a L_a^s L_a^t ,
\label{eq:ch2V}
\end{equation}
%
and on a favourable CICY effectiveness of a curve class is componentwise
non-negativity in the basis descending from the ambient factors.

It is tempting to conclude that \cref{eq:ch2V} makes the condition automatic:
the triple intersection numbers satisfy $d_{rst} \ge 0$, and
$\sum_a L_a L_a$ is positive semi-definite as a matrix. The inference fails,
and the reason is a useful one to record. A positive semi-definite matrix has
negative \emph{off-diagonal} entries --- $L = (1,-1,0,0)$ is the smallest
example --- so the contraction
%
\begin{equation}
\underbrace{d_{rst}}_{\ge\,0}
\underbrace{\textstyle\sum_a L_a^s L_a^t}_{\text{positive semi-definite}}
\qquad\text{is \emph{not} sign definite,}
\end{equation}
%
and the anomaly genuinely constrains. 

On the tetraquadric the charges
$\{(3,-3,0,0), (-3,3,0,0), 0,0,0\}$ 
\newline give surplus $(24,24,-12,-12)$ and are excluded; the same shape at charge two is admitted.

\subsection{Poly-stability, and an exact sign obstruction}
\label{sec:slopeobstruction}

The Hermitian--Yang--Mills equation has a solution when $V$ is poly-stable,
which for \cref{eq:lbsum} means every summand has vanishing slope
simultaneously,
%
\begin{equation}
\mu(L_a)(t) \;=\; d_{rst}\, L_a^r t^s t^t \;=\; 0
\quad\text{for all } a,
\qquad t \in \text{K\"ahler cone} .
\label{eq:slope}
\end{equation}
%
Since $\sum_a L_a = 0$ the slopes sum to zero identically, so \cref{eq:slope}
is $n-1$ conditions on the $h^{1,1}-1$ projective directions of the cone: a
common zero of several quadrics, generically empty, and expensive to search
for numerically.

There is an exact obstruction which is pure sign arithmetic. Write
$M_a^{st} = d_{rst} L_a^r$, symmetrised, so that $\mu(L_a) = t^{\top} M_a t$.
At a common zero $t_\star$ of \emph{all} the slopes, every \emph{partial} sum
vanishes as well:
%
\begin{equation}
\sum_{a \in S} \mu(L_a)(t_\star)
\;=\; t_\star^{\top}
\underbrace{\Bigl(\sum_{a \in S} M_a\Bigr)}_{=:\,M_S} t_\star
\;=\; 0
\qquad\text{for every subset } S .
\label{eq:subsetobstruction}
\end{equation}
%
On a favourable CICY the K\"ahler cone is the positive orthant, so if $M_S$ is
non-zero with no negative entry then $t^\top M_S t > 0$ throughout the interior
and no such $t_\star$ exists. Hence:
%
\begin{equation}
\exists\, S \subseteq \{1,\dots,n\} :\;
M_S \neq 0 \ \text{and} \ M_S \ \text{has no negative entry}
\quad\Longrightarrow\quad
V \ \text{is nowhere poly-stable.}
\label{eq:obstruction}
\end{equation}
%
\Cref{eq:obstruction} is $2^n - 2$ sign tests --- at most thirty for the ranks
of interest --- against a numerical search costing some thirty milliseconds.
Measured on a sample of twenty thousand rank-five sums on the tetraquadric it
removed every one that the numerical search would have removed, at
$3.5\times10^{-2}$ milliseconds each. It is also \emph{inherited}: if a
subset already fails, so does every completion of it, which is what makes an
exhaustive scan finite. The subsets of size one are the statement that a
summand whose slope form is one-signed can never appear in a poly-stable
bundle, and this survives as a condition on individual line bundles that can
be applied to a candidate pool before any sum is assembled.

One boundary case matters. If $M_a = 0$ the slope vanishes \emph{identically},
which is compatible with poly-stability rather than an obstruction to it; the
trivial summand $\cO_X$ is the obvious instance. Treating $M_a = 0$ as
``not sign-changing'' silently deletes valid models.

\subsection{A parity obstruction on the tetraquadric}

On the tetraquadric $\bigl[\PP^1\,2;\PP^1\,2;\PP^1\,2;\PP^1\,2\bigr]$ the only
non-zero entry of $d_{rst}$ is $2$, so $6\operatorname{ind}(V)$ is even for
every line bundle sum and
%
\begin{equation}
\operatorname{ind}(V) \;=\; -3
\qquad\text{is unreachable, at any charge.}
\label{eq:tetraparity}
\end{equation}
%
Three generations on that manifold therefore \emph{requires} a non-trivial
freely acting quotient --- the same conclusion \cref{sec:generations} reaches
for the standard embedding by the entirely different route of
\cref{eq:genquot}, and a small illustration that the two constructions are
constrained by the same arithmetic.

% =====================================================================
\section{Equivariant structures and the character-valued index}
\label{sec:equivariant}

\Cref{eq:genquot} divides a generation count by $|\Gamma|$. That is arithmetic,
not a model: it presumes the bundle descends to $X/\Gamma$, which requires an
\emph{equivariant structure} --- a lift of the $\Gamma$ action to the total
space of each line bundle. The structure is extra data again, and different
lifts give different spectra on identical topological input.

\subsection{What is exactly computable, and why}

The individual groups $H^q(X, L)$ as $\Gamma$-modules require the Leray
spectral sequence run equivariantly, and the ranks of its differentials are
not fixed by the degrees. The \emph{alternating sum} is different, because
characters are additive on exact sequences. Any equivariant resolution
therefore computes
%
\begin{equation}
\operatorname{ind}_\Gamma(L)
\;=\; \sum_q (-1)^q \operatorname{ch} H^q(X, L) \;\in\; R(\Gamma) ,
\end{equation}
%
with no spectral sequence anywhere. The Koszul resolution of $\cO_X$ inside
the ambient product,
%
\begin{equation}
0 \to \Lambda^K N^* \to \cdots \to N^* \to \cO_A \to \cO_X \to 0,
\qquad N^* = \bigoplus_a \cO_A(-d_a),
\end{equation}
%
gives every term as a sum of ambient line bundles, whose equivariant Euler
characteristics factorise over the projective factors by K\"unneth. Writing
$c_a$ for the character by which $\Gamma$ acts on the $a$-th defining
polynomial,
%
\begin{equation}
\operatorname{ind}_\Gamma\bigl(\cO_X(k)\bigr)
\;=\;
\sum_{S \subseteq \{1,\dots,K\}} (-1)^{|S|}
\underbrace{\Bigl(\prod_{a \in S} c_a^{-1}\Bigr)}_{\text{Koszul twist}}
\;
\underbrace{\chi_\Gamma\Bigl(A,\ \cO_A\bigl(k - \textstyle\sum_{a\in S} d_a\bigr)\Bigr)}_{\text{K\"unneth product of monomial characters}} ,
\label{eq:koszulindex}
\end{equation}
%
a finite sum in exact integer arithmetic. Its total must agree with the
ordinary index of \cref{eq:lineindex}, and does; the Koszul route is in fact
the more exact of the two, returning integers where the spectral-sequence
route returns floating point.

When $\Gamma$ permutes the ambient factors, \cref{eq:koszulindex} still holds
with one change: $g^j$ permutes the tensor factors, so its trace is a product
over the \emph{cycles} of $\sigma^j$ rather than over the factors, the
contribution of a cycle of length $\ell$ through $i$ being the trace of the
composite once around, diagonal with weights $\sum_{s<j\ell} w_{\sigma^s(i)}$.
When $\Gamma$ permutes the defining polynomials, only subsets with
$\pi^j(S) = S$ contribute, each carrying the sign of $\pi^j$ restricted to $S$
--- the cost of reordering $e_{\pi(a_1)} \wedge \cdots \wedge e_{\pi(a_r)}$
back into increasing order.

\begin{remark}[A sign the usual check cannot see]
That wedge sign is invisible to the comparison against the ordinary index:
at the identity $\pi^0$ is trivial and the sign is $+1$, so forcing it to
$+1$ everywhere still reproduces the total exactly. It needs a separate
oracle. One is available: on a configuration whose two defining polynomials
have the same multidegree --- $[\PP^1\,1\,1]^{\times 5}$, a favourable
threefold with $\chi = -80$ --- swapping them is, in the eigenbasis
$p_\pm = p_1 \pm p_2$, precisely the phase-only action with charges $0$ and
$1$. The two computations agree; forcing the sign wrong breaks $16$ of $25$
bundles while leaving the total untouched.
\end{remark}

\subsection{Freeness, twice}

If $\Gamma$ acts freely then every $g \neq e$ has no fixed points, so its
holomorphic Lefschetz number vanishes, and
%
\begin{equation}
L(g) = 0 \ \ \forall g \neq e
\qquad\Longleftrightarrow\qquad
\operatorname{ind}_\Gamma(L) \ \text{is a multiple of the regular
representation,}
\label{eq:lefschetzfree}
\end{equation}
%
that is, the charges are equidistributed. \Cref{eq:lefschetzfree} is necessary
and not sufficient, and it depends on the set of bundles probed: an action
trivial on two of the four tetraquadric factors --- plainly not free, with a
two-dimensional fixed locus --- passes a hand-picked list of five probes and
fails $180$ of a $625$-probe box.

A second, independent test comes from the geometry. A diagonal action fixes
the coordinate points of the ambient, and at such a point every monomial
vanishes except the pure power of the chosen coordinates; if that monomial's
charge differs from the declared charge of a defining polynomial, the
polynomial vanishes there and the fixed point \emph{lies on} $X$. Neither test
is complete --- the geometric one sees only coordinate fixed points, the
Lefschetz one is necessary only --- and they agree wherever both can see.

\begin{remark}[A case the hand argument gets wrong]
On the tetraquadric the $\mathbb{Z}_2$ action $[x_0:x_1] \mapsto [x_0:-x_1]$
on each factor is free: the sixteen ambient fixed points carry monomials of
charge $8 \equiv 0 \bmod 2$, which may appear in a charge-zero quartic, so a
generic such quartic misses them. The identical weights for $\mathbb{Z}_3$,
$\mathbb{Z}_4$ and $\mathbb{Z}_5$ are \emph{not} free: $11$, $8$ and $15$ of
the sixteen points are forced onto $X$, because $8 \not\equiv 0$ modulo those
orders. The naive argument ``a generic hypersurface avoids finitely many
points'' would accept all four.
\end{remark}

\subsection{The structure does not matter, for the chiral spectrum}
\label{sec:structureirrelevant}

The equivariant structures on a line bundle form a torsor over the character
group of $\Gamma$: for cyclic $\Gamma$ of order $n$ there are exactly $n$ of
them, and a rank-$r$ sum admits $n^r$, of which those with trivial total twist
descend. This is a finite but real ambiguity.

For a \emph{free} action it is inert. By \cref{eq:lefschetzfree} each summand's
index character is a multiple of the regular representation, hence a constant
vector; a twist permutes a constant vector cyclically; a constant vector is
fixed by any permutation. Therefore
%
\begin{equation}
\Gamma \ \text{free}
\quad\Longrightarrow\quad
\underbrace{\operatorname{ind}_\Gamma(V)}_{\text{chiral spectrum}}
\ \text{is independent of which of the } n^r \ \text{structures is chosen.}
\label{eq:structureirrelevant}
\end{equation}
%
The choice shows up only in $h^1$ and $h^2$ separately --- the vector-like
pairs --- which an index cannot see in any case. So the ambiguity is real and
confined to exactly the sector where nothing here could have computed anything
anyway.

\subsection{Which lattice, and which group}

Two bookkeeping points are easy to get wrong and change the answer.

\emph{The lattice.} Reflexivity, and hence the whole Batyrev construction of
\cref{sec:polytopes}, is a property of a polytope \emph{and a lattice}. The
$24$-cell written as the $D_4$ roots $\pm e_i \pm e_j$ has polar dual with
vertices $\{\pm e_i\} \cup \{\tfrac12(\pm1,\pm1,\pm1,\pm1)\}$, which is not
integral in $\mathbb{Z}^4$: the polytope is \emph{not} reflexive there. It is
reflexive with respect to the lattice its own vertices generate,
$D_4 = \{x \in \mathbb{Z}^4 : \sum x_i \in 2\mathbb{Z}\}$, of index two.

\emph{The group.} An element acts on the ambient by a permutation of factors
and linear maps, and two such data give the same map exactly when they differ
by a scalar on each factor. The defining polynomials do not respect that
quotient: rescaling factor $i$ by $\lambda_i$ multiplies $p_a$ by
$\prod_i \lambda_i^{d[i][a]}$. Element equality is projective on the ambient
and \emph{not} projective on the polynomials, and one cannot have both. Working
with an explicit lift --- the group generated under composition with exact
equality on the full data --- makes every character unambiguous, at the cost
that $|\Gamma|$ then refers to the lift. If the lift contains a non-identity
element acting trivially on the ambient it is a proper central extension, and
the order to divide by in \cref{eq:genquot} is the geometric one. Rescaling one
tetraquadric factor by $-1$ is the smallest example: trivial on $X$, of order
two on $\cO(k)$, and at $k = (1,0,0,0)$ none of the two sections descend while
at $k = (2,0,0,0)$ all four do.

% =====================================================================
\section{Wilson lines and the Standard Model group}
\label{sec:breaking}

A rank-five $V$ leaves $SU(5)$ unbroken in four dimensions. Breaking it needs a
Wilson line on $X/\Gamma$, and which groups can supply one is a question of
pure arithmetic.

\subsection{Branching}

With the hypercharge generator $Y = \operatorname{diag}(-\tfrac13,-\tfrac13,
-\tfrac13,\tfrac12,\tfrac12)$ in the fundamental,
%
\begin{align}
\mathbf{10} \;&\longrightarrow\;
(\mathbf{3},\mathbf{2})_{1/6}
\;+\; (\overline{\mathbf{3}},\mathbf{1})_{-2/3}
\;+\; (\mathbf{1},\mathbf{1})_{1} ,
\label{eq:ten}\\
\overline{\mathbf{5}} \;&\longrightarrow\;
(\overline{\mathbf{3}},\mathbf{1})_{1/3}
\;+\; (\mathbf{1},\mathbf{2})_{-1/2} ,
\label{eq:fivebar}
\end{align}
%
one Standard Model generation in the left-handed conjugate convention, with
$\operatorname{Tr} Y = 0$ over $\mathbf{10} + \overline{\mathbf{5}}$.

\subsection{Which cyclic groups can break $SU(5)$}

A Wilson line in the hypercharge direction is
$W = \operatorname{diag}(a,a,a,b,b)$ with $\det W = a^3 b^2 = 1$. Its
commutant in $SU(5)$ is $S\bigl(U(3)\times U(2)\bigr) = SU(3)\times SU(2)\times
U(1)$ when $a \neq b$, and all of $SU(5)$ when $a = b$, in which case $W$ is
central and breaks nothing. Writing $a = \omega^p$, $b = \omega^q$ with
$\omega = e^{2\pi i/n}$, the determinant condition is $3p + 2q \equiv 0
\pmod n$, and
%
\begin{equation}
\#\bigl\{(p,q) : 3p + 2q \equiv 0 \ (n),\ p \neq q\bigr\}
\;=\;
\underbrace{n}_{\substack{\text{kernel of }(p,q)\mapsto 3p+2q,\\ \text{surjective since }\gcd(2,3,n)=1}}
\;-\;
\underbrace{\gcd(n,5)}_{\substack{\text{those with } p=q,\\ \text{i.e. } 5p \equiv 0}} .
\label{eq:wilsoncount}
\end{equation}
%
\Cref{eq:wilsoncount} vanishes exactly when $n \mid 5$, that is for $n = 1$ and
$n = 5$. So
%
\begin{equation}
\Gamma = \mathbb{Z}_5
\quad\text{divides the generation count by five and \emph{cannot} break }SU(5),
\end{equation}
%
every compatible Wilson line being central, while $\mathbb{Z}_2$ is the
smallest group that can --- and is exactly the order that
\cref{eq:tetraparity} forces on the tetraquadric. The two constraints, one from
the parity of the intersection numbers and one from the determinant of a
Wilson line, land on the same group for unrelated reasons.

\subsection{The chiral spectrum descends without splitting}

The $\mathbf{10}$ sits in $H^1(V)$ and the $\overline{\mathbf{5}}$ in
$H^1(\Lambda^2 V)$, so their net counts are read off
$\operatorname{ind}(V)$ and $\operatorname{ind}(\Lambda^2 V)$. A Wilson line
shifts the $\Gamma$-charge of each Standard Model piece within its multiplet by
a \emph{different} amount --- that is how it splits them --- and the surviving
multiplicity is the one at the shifted charge.

But by \cref{eq:structureirrelevant} every index character of a free action is
constant, so every shift lands on the same multiplicity:
%
\begin{footnotesize}
\begin{equation}
\underbrace{\text{Wilson line}}_{\text{shifts each piece differently}}
\ \text{acting on} \
\underbrace{\text{a constant character}}_{\text{\cref{eq:lefschetzfree}}}
\quad\Longrightarrow\quad
\text{complete } SU(5) \ \text{generations, } \
-\operatorname{ind}(V)/|\Gamma| \ \text{of them.}
\label{eq:nosplit}
\end{equation}
\end{footnotesize}
%
The Wilson line cannot split the chiral spectrum at all. This is the correct
physics and it resolves an apparent tension: doublet--triplet splitting is
necessarily a statement about \emph{vector-like} pairs, the non-chiral content
that an index cannot see, which is precisely the half left undetermined by
\cref{eq:structureirrelevant}.

For other ranks the commutant differs --- $SU(4) \to SO(10)$, $SU(3) \to E_6$
--- and two identities serve as free checks: for rank four $\Lambda^2 V$ is the
self-dual $\mathbf{6}$ of $SU(4)$, so $\operatorname{ind}(\Lambda^2 V) = 0$
identically and the $\mathbf{10}$ of $SO(10)$ is vector-like; for rank three
$\Lambda^2 V \cong V^*$, so $\operatorname{ind}(\Lambda^2 V) =
-\operatorname{ind}(V)$ identically. Both hold at the level of the whole
$\Gamma$-character, not merely the total.

% =====================================================================
\section{Yukawa couplings: what is exact}
\label{sec:yukawa}

\Cref{sec:notdetermined} says masses and couplings are not functions of the
configuration matrix. That is true of \emph{physical} couplings and needs
qualifying, because one holomorphic coupling is exactly computable and the
entire \emph{pattern} of another is.

\subsection{The standard embedding: an integer series}

For $V = TX$ the coupling of three $(1,1)$-type families is
%
\begin{equation}
Y_{rst} \;=\; \int_X J_r \wedge J_s \wedge J_t \;=\; d_{rst} ,
\label{eq:classicalyukawa}
\end{equation}
%
the triple intersection numbers of \cref{sec:cipro} --- integers, no metric
anywhere. \Cref{eq:classicalyukawa} is the classical limit; the coupling
receives worldsheet instanton corrections controlled by the genus-zero
Gopakumar--Vafa invariants of \cref{sec:gv},
%
\begin{equation}
Y(q) \;=\;
\underbrace{d_{rst}}_{\text{\cref{eq:classicalyukawa}}}
\;+\;
\underbrace{\sum_{\beta} n_\beta\, \beta_r \beta_s \beta_t\,
\frac{q^{\beta}}{1 - q^{\beta}}}_{\text{worldsheet instantons}} ,
\label{eq:quantumyukawa}
\end{equation}
%
so that on the quintic
%
\begin{equation}
Y(q) \;=\; 5 \;+\; 2875\,\frac{q}{1-q}
\;+\; 8 \cdot 609250\, \frac{q^2}{1-q^2}
\;+\; 27 \cdot 317206375\, \frac{q^3}{1-q^3} \;+\; \cdots ,
\end{equation}
%
every coefficient an integer. This is a genuine Yukawa coupling, computed
exactly, and the same enumerative machinery that verifies the mirror map in
\cref{sec:gv} produces it.

\begin{remark}[The series converges only for small $q$]
The invariants grow faster than $q^d$ shrinks unless $q$ is genuinely small.
On the quintic $n_5 = 2.3\times10^{14}$, so at $q = 10^{-2}$ the degree-five
term alone is $3\times10^{6}$ and the partial sum is meaningless; at
$q = 10^{-4}$ the last term is $3\times10^{-4}$ and the sum is a coupling.
A truncation of \cref{eq:quantumyukawa} should always be reported with the
size of its last term.
\end{remark}

\subsection{Line bundle models: the texture is exact, the values are not}

For \cref{eq:lbsum} of rank five the matter is
%
\begin{equation}
\mathbf{10}_a \leftrightarrow H^1(X, L_a),
\qquad
\overline{\mathbf{5}}_{ab} \leftrightarrow H^1(X, L_a \otimes L_b),
\qquad
\mathbf{5}_{ab} \leftrightarrow H^1(X, L_a^{-1} \otimes L_b^{-1}),
\end{equation}
%
and a coupling is a cup product landing in
$H^3(X, \Lambda^3 V) $. That is $H^3(\cO_X) = \mathbb{C}$ only when the
participating line bundles cancel exactly, which gives a selection rule
needing no representatives at all:
%
\begin{align}
\text{up-type} \quad
&\mathbf{10}_a\, \mathbf{10}_b\, \mathbf{5}_{ab}
&&\text{allowed for every } a < b,
\label{eq:upsel}\\[2pt]
\text{down-type} \quad
&\mathbf{10}_a\, \overline{\mathbf{5}}_{bc}\, \overline{\mathbf{5}}_{de}
&&\text{allowed only when } \{a,b,c,d,e\} \ \text{are all distinct.}
\label{eq:downsel}
\end{align}
%
\Cref{eq:downsel} is worth pausing on. Since $\sum_a L_a = 0$, the charges
cancel precisely when all five indices appear, so
%
\begin{equation}
\underbrace{\varepsilon_{abcde}}_{\text{the } SU(5) \ \text{invariant}}
\qquad\text{reappears as}\qquad
\underbrace{\textstyle\sum_a L_a = 0}_{\text{the } S(U(1)^5) \ \text{condition}} ,
\end{equation}
%
and there are exactly $5 \times 3 = 15$ down-type patterns: five choices of
$a$, three ways of pairing the remaining four.

A second filter is equally exact. A charge-allowed coupling is still absent if
any of the three cohomology groups is zero-dimensional --- there is no field to
couple --- and those dimensions are computed by the same line bundle
cohomology that gives the spectrum. The result is a \emph{texture} of allowed
and forbidden couplings, forbidden in two distinct senses:
%
\begin{equation}
\underbrace{\text{charge-forbidden}}_{\text{symmetry}}
\qquad\text{versus}\qquad
\underbrace{\text{charge-allowed but } h^1 = 0}_{\text{geometry: a \emph{texture zero}}} .
\end{equation}
%
What is \emph{not} available is the value of a surviving coupling. It is
quasi-topological --- no metric required --- but evaluating the cup product
needs explicit cohomology representatives rather than dimensions. That is a
missing computation, not an obstruction, and it should be kept distinct from
the physical coupling, which \emph{is} obstructed.

\subsection{A necessary condition on the manifold}

Since an up-type coupling needs $h^1(L_a) > 0$, $h^1(L_b) > 0$ and
$h^1(L_a^{-1} L_b^{-1}) > 0$ with charges summing to zero, one may ask whether
any such triple exists at all inside a charge box --- a question about the
manifold, answerable before any bundle is chosen. If none does, no model in
that box has a coupling, whatever else it satisfies.

The condition sharpens when combined with \cref{sec:slopeobstruction}: a
summand whose slope form is one-signed can never sit in a poly-stable bundle,
so each member of the triple must also be slope-non-definite. The two filters
together are informative. Without the slope condition, the first triple found
on the tetraquadric contains $(0,0,-2,0)$, whose slope is definite; a search
over $9390$ models built around that pair contains exactly zero poly-stable
ones. With both conditions, viable triples exist on the tetraquadric at charge
two and on CICY~$7833$ at charge five but not at charge four --- a reminder
that the answer is a statement about the box as much as the manifold. Both
conditions are necessary and neither is sufficient: a model realising a viable
triple must still close with the right index, anomaly and \emph{joint}
poly-stability, and the per-summand slope test does not imply the joint one.

% =====================================================================
\section{Reflexive polytopes and the Batyrev construction}
\label{sec:polytopes}

The toric side of the same package is two-dimensional in
\cite{SuppFigures}: sixteen reflexive polygons, each the toric diagram of a
local Calabi--Yau. Reflexivity generalises, and in four dimensions it produces
compact threefolds.

A lattice polytope $P$ containing the origin in its interior is
\emph{reflexive} when its polar dual
$P^* = \{y : \langle x,y\rangle \ge -1 \ \forall x \in P\}$ is again a lattice
polytope. For a reflexive four-polytope $\Delta$ the anticanonical
hypersurface in the toric variety of the face fan of $\Delta^*$ is a
Calabi--Yau threefold with
%
\begin{align}
h^{1,1} &= \ell(\Delta^*) - 5
- \sum_{\theta^* \ \mathrm{facet}} \ell^*(\theta^*)
+ \sum_{\theta^* \ \mathrm{codim}\,2} \ell^*(\theta^*)\, \ell^*(\theta) ,
\label{eq:batyrev11}\\
h^{2,1} &= \ell(\Delta) - 5
- \sum_{\theta \ \mathrm{facet}} \ell^*(\theta)
+ \sum_{\theta \ \mathrm{codim}\,2} \ell^*(\theta)\, \ell^*(\theta^*) ,
\label{eq:batyrev21}
\end{align}
%
with $\ell$ the number of lattice points of a face and $\ell^*$ the number in
its relative interior, and $\theta \leftrightarrow \theta^*$ the
dimension-reversing duality $\dim\theta + \dim\theta^* = 3$.

\subsection{The 24-cell}

The $24$-cell is reflexive in $D_4$ but not in $\mathbb{Z}^4$, as noted in
\cref{sec:equivariant}. In a basis of $D_4$ both it and its dual are lattice
polytopes with $f$-vector $(24, 96, 96, 24)$, and
%
\begin{equation}
\ell(\Delta) \;=\; \ell(\Delta^*) \;=\;
\underbrace{24}_{\text{vertices}} + \underbrace{1}_{\text{origin}} \;=\; 25 ,
\end{equation}
%
so no proper face has an interior lattice point and \emph{both} correction sums
in \cref{eq:batyrev11,eq:batyrev21} vanish identically. Hence
%
\begin{equation}
h^{1,1} \;=\; h^{2,1} \;=\; 25 - 5 \;=\; 20 ,
\qquad \chi \;=\; 0 ,
\end{equation}
%
and self-duality forces $h^{1,1} = h^{2,1}$ independently of the arithmetic, so
two arguments land in the same place. Braun's Hodge numbers
$(1,1)$~\cite{Braun2011} are those of a free quotient of this threefold, not
of the threefold itself.

\begin{remark}[Which polytope to pass]
$\Delta$ in \cref{eq:batyrev11,eq:batyrev21} is the \emph{Newton} polytope,
whose lattice points index monomials, and $\Delta^*$ the fan polytope, whose
vertices are the rays. Interchanging them transposes the Hodge numbers, and
the failure is silent. The quintic calibrates it: the fan polytope of $\PP^4$
has six lattice points and its polar has $\binom{9}{4} = 126$, one per quintic
monomial, giving $(h^{1,1}, h^{2,1}) = (1, 101)$; passing the fan polytope
directly returns $(101,1)$, the mirror. For a self-dual polytope such as the
$24$-cell the distinction is invisible, which is exactly why it needs stating.
\end{remark}

% =====================================================================
\section{The characteristic polynomial of the Hofstadter model}
\label{sec:hofstadter}

The quantized mirror curves of \cite{SuppFigures} produce Hofstadter spectra
numerically. There is a closed form for the characteristic polynomial, due to
Marra, Proietti and Sheng~\cite{MarraProiettiSheng2024}, and it connects to the
same enumerative geometry through the relativistic Toda lattice.

Let $\tilde e_k$ denote the \emph{two-step} elementary symmetric polynomial,
summing over $k$-subsets no two of whose indices are adjacent:
%
\begin{equation}
\tilde e_k(x_1,\dots,x_m)
\;=\;
\sum_{\substack{1 \le j_1 < \cdots < j_k \le m \\ |j_i - j_{i+1}| \ge 2}}
x_{j_1} \cdots x_{j_k} ,
\qquad \tilde e_0 = 1 .
\end{equation}
%
These are what appear in the determinant of a tridiagonal matrix, because a
term of the Leibniz expansion can use each off-diagonal pair at most once and
never two overlapping ones. At the mid-band point the characteristic
polynomial of the flux-$P/Q$ Hofstadter Hamiltonian is
%
\begin{equation}
f(E) \;=\; \det\bigl(\hat H_{P/Q} - E\bigr)
\;=\;
\sum_{i=0}^{\lfloor Q/2 \rfloor} (-1)^{Q+i}\, 4^i\, E^{Q-2i}\,
\tilde e_i\bigl(\sin^2(\pi\alpha), \dots, \sin^2((Q-1)\pi\alpha)\bigr),
\label{eq:hofstadterpoly}
\end{equation}
%
with $\alpha = P/Q$. Only powers $E^{Q-2i}$ survive, which is the statement
that the spectrum is symmetric under $E \mapsto -E$ for even $Q$.

Two consequences are worth recording. The constant term gives a closed-form
evaluation with no free parameters and no $P$ dependence,
%
\begin{equation}
\tilde e_{Q/2}\bigl(\sin^2(\pi\alpha), \dots, \sin^2((Q-1)\pi\alpha)\bigr)
\;=\; 4^{-(Q/2 - 1)}
\qquad (Q \ \text{even}),
\label{eq:etildeidentity}
\end{equation}
%
and the dependence on the Brillouin torus is a single additive constant,
independent of $E$:
%
\begin{equation}
f(E, \nu_x, \nu_y) \;=\;
f\Bigl(E, \tfrac{\pi}{2Q}, \tfrac{\pi}{2Q}\Bigr)
\;+\; 2(-1)^{Q-1}\bigl(\cos Q\nu_x + \cos Q\nu_y\bigr) ,
\label{eq:chambers}
\end{equation}
%
the Chambers relation. \Cref{eq:chambers} is why the band structure is a rigid
translation of one polynomial rather than a family of unrelated ones, and it
fixes where the zero mode sits: $E = 0$ occurs at the mid-band point when $Q$ is
odd, at the corner when $Q \equiv 2 \bmod 4$, and at the centre when
$4 \mid Q$.

\begin{remark}[Computability]
$\tilde e_k$ read as a definition is exponential in the number of variables.
The recurrence $E_k(m) = E_k(m-1) + x_m E_{k-1}(m-2)$ --- a subset either omits
$x_m$, or contains it and then cannot contain $x_{m-1}$ --- makes it quadratic,
which is the difference between $2^{200}$ subsets and three milliseconds for
the degree-$201$ polynomial. That is the range in which the butterfly is
actually drawn.
\end{remark}

% =====================================================================
\section{The boundary, stated precisely}
\label{sec:boundary}

\Cref{sec:notdetermined} draws the line in one sentence. It is worth drawing it
in three, because the failures are different in kind and conflating them makes
a construction look more predictive than it is.

\begin{enumerate}[leftmargin=2em]
\item \emph{Not implemented, but exact in principle.} The values of the
surviving line bundle Yukawa couplings of \cref{sec:yukawa}. They are
quasi-topological; they need explicit cohomology representatives, obtained
from a Koszul or \v{C}ech resolution, and no metric. Supplying them would give
exact integers, as \cref{eq:classicalyukawa} does for the standard embedding.

\item \emph{Obstructed by the metric.} The physical couplings. A physical
Yukawa is the holomorphic one divided by the norms of the three fields, and
those norms are integrals of harmonic representatives against the Ricci-flat
metric, for which no closed form is known on any compact Calabi--Yau
threefold. Numerical approximations exist~\cite{Larfors2022,HendiLarforsWalden2024, Butbaia2024}, and
handing them a model requires supplying, beyond the polytope data, three
things a metric package cannot know: that the defining polynomial must be
$\Gamma$-invariant, that the K\"ahler moduli must sit where the bundle is
poly-stable, and that $\Gamma$ must preserve the holomorphic form. The last is
a condition on characters,
%
\begin{equation}
\chi(\Omega) \;=\;
\underbrace{\sum_{\text{all coordinates}} w}_{\text{ambient volume form}}
\;-\;
\underbrace{\sum_a c_a}_{\text{Jacobian of the defining polynomials}}
\;\equiv\; 0 ,
\label{eq:omegachar}
\end{equation}
%
which is $\Gamma \subset SU(3)$ rather than merely $U(3)$; if
\cref{eq:omegachar} fails then $\Omega$ does not descend and $X/\Gamma$ is not
Calabi--Yau at all. It is independent of freeness and neither implies the
other: on the tetraquadric the order-four action preserves $\Omega$ and is not
free, while the $\mathbb{Z}_3$ action does neither.

\item \emph{Obstructed by unsolved physics.} The masses. A fermion mass is a
physical coupling times a Higgs vacuum expectation value, and the vacuum
expectation value requires the moduli to be stabilised.
\end{enumerate}

\noindent
Nothing in \cref{sec:bundles,sec:equivariant,sec:breaking,sec:yukawa} touches
the second or third. What they do is push the first boundary as far as it will
go, and mark exactly where it stops.

% =====================================================================

\section{Conclusion: transport without a carrier}
\label{sec:aside}
What follows is an analogy, not a result.  The Heterotic anomaly cancellation reads
%
\begin{equation}
\ch_2(T_X) \;=\; \ch_2(V) \;-\; [C],
\label{eq:anomaly}
\end{equation}
%
with $V$ the gauge bundle and $[C]$ the class of curves wrapped by
five-branes~\cite{Anderson2025,Anderson2023branes}. Across a conifold
transition the left-hand side changes,
%
\begin{equation}
\ch_2(T_{\XR}) \;=\; \ch_2(T_{\XD}) \;+\; \bigl[\PP^1_s\bigr],
\qquad\text{yet \cref{eq:anomaly} must hold on both sides,}
\label{eq:ch2}
\end{equation}
%
and the resolution offered in~\cite{Anderson2025} is that the bridging curve
classes shift in compensation,
%
\begin{equation}
[C_R] \;=\; [C_D] \;-\; \bigl[\PP^1_s\bigr].
\label{eq:bridging}
\end{equation}
%
The conserved quantity is neither destroyed nor independently created: it changes
which object carries it. The same pattern appears in the moduli map, where
active complex structure moduli on the deformation side correspond to brane
and bundle moduli together with the new K\"ahler modulus on the resolution
side, while computable quantities are unchanged --- in one worked example
of~\cite{Anderson2025}, $147{,}440$ independent Yukawa couplings agree across
the correspondence. The labels move; the observables do not.

Aharonov, Collins and Popescu~\cite{Aharonov2023} exhibit angular momentum
transported across a region where the probability of finding any particle or
field is vanishingly small, and argue that the picture of a conservation law
as a flux of carriers needs revisiting. Waegell, Tollaksen and
Aharonov~\cite{Waegell2024} discuss a pre- and post-selected arrangement in
which, on their reading, a particle's mass and momentum are localised in
different places. In both cases, as in \cref{eq:ch2,eq:bridging}, a conserved
bookkeeping quantity is not tied to the object one would naively call its
carrier.

The resemblance should not be pushed. \Cref{eq:ch2} is an identity between
cohomology classes on distinct manifolds, whereas weak values are complex
numbers extracted from an ensemble. Nothing in
\cref{eq:anomaly,eq:ch2,eq:bridging} describes a process in time, while
\cite{Aharonov2023} is precisely about flow through an intermediate region.
The quantum reading is contested: the authors of~\cite{Waegell2024} state in
their own abstract that the titular interpretation is extremely controversial
and rests on several assumptions, and offer an alternative. And there is no
map from one setting to the other; the phrase ``a conserved quantity not
residing where one expects'' is doing all of the work and is broad enough to
cover much else. The honest summary is that both literatures make the same
pedagogical point --- identifying a conserved quantity with a carrier is a
habit rather than a necessity --- which is a reason to find both interesting,
not evidence that they are the same phenomenon.

% =====================================================================
\appendix
\small

\section{Source Code and Limitations}
\label{sec:limits}

\begin{enumerate}[leftmargin=*,itemsep=2pt]

\item Source code: \url{https://github.com/brentharts/CICY}

\item \textbf{The web is partial.} Depth $\FactDepth$ yields $\FactConfigs$
configurations and $\FactHodgePairs$ distinct Hodge pairs, against $7890$ and
$266$ in the published list, with the frontier still expanding. Going deeper
is a compute budget question.

\item \textbf{$\ch_2$ is checked only in the split presentation}, where
\cref{eq:reembed} puts both sides in one ambient space. Transitions requiring
a general toric weight matrix, such as the flop partner of Anderson and Gray (2025) ~\cite{Anderson2025}
eq.~(1.14), are outside what a configuration matrix expresses. The ambient
class representation is also not faithful, since a class can restrict to zero
on $X$.

\item \textbf{Gopakumar--Vafa invariants cover only the one-parameter models}
of \cref{eq:oneparam}, at genus zero. Every other CICY has $h^{1,1}>1$ and
needs the multivariable mirror map, which \texttt{CIPro} Anderson and Constantin (2026) ~\cite{Anderson2026}
covers and this does not.

\item \textbf{Fourfold Hodge numbers are lightly tested.} One case, the
sextic in $\PP^5$, is verified here; the four-fold
classification~\cite{GrayHauptLukas2013} is not otherwise exercised.

\item \textbf{Smoothness is checked over finite fields, not proved.} The
Jacobian rank criterion of~\cite{Anderson2026} is implemented by exhaustive
enumeration over $\PP^{n_1}(\mathbb{F}_p)\times\cdots$, which is rigorous
about the reduction modulo $p$ but is evidence, not proof, in characteristic
zero; the implementation returns that distinction rather than a bare boolean.
Two traps are handled: a prime dividing a degree makes formal derivatives
vanish so that everything looks singular, and a single random member over a
small field need not be generic, so the generic question is settled by
exhibiting one smooth member. The published singular
example of Anderson, et al. (2026) ~\cite{Anderson2026} is correctly reported as singular.

\item \textbf{Only the standard embedding is treated.} \Cref{eq:gen} assumes
$V = TX$. Realistic model building uses other bundles~\cite{Anderson2011,%
Anderson2012}, for which the spectrum depends on $V$ and on Wilson lines and
not on the configuration matrix alone. Nothing here yet computes a mass or a
coupling; see \cref{sec:notdetermined}.
\end{enumerate}



% =====================================================================
\small
\begin{thebibliography}{99}
\setlength{\itemsep}{1pt plus 0.5pt}

\bibitem{Anderson2026}
L.~B.~Anderson, A.~Constantin, J.~Gray, Y.-H.~He, S.-J.~Lee and A.~Lukas,
\emph{CIPro Package: Complete Intersections in Products of Projective Spaces
and Line Bundles},
\href{https://arxiv.org/abs/2606.27588}{arXiv:2606.27588}.


\bibitem{Candelas1988}
P.~Candelas, A.~M.~Dale, C.~A.~L\"utken and R.~Schimmrigk,
\emph{Complete Intersection Calabi--Yau Manifolds},
Nucl.\ Phys.\ B \textbf{298} (1988) 493.

\bibitem{Anderson2025}
L.~B.~Anderson, J.~Gray, S.~A.~Patil and C.~Scanlon,
\emph{Mapping moduli across heterotic conifolds},
\href{https://arxiv.org/abs/2512.18124}{arXiv:2512.18124}.

\bibitem{Larfors2019}
M.~Larfors and R.~Schneider,
\emph{Line bundle cohomologies on CICYs with Picard number two},
Fortsch.\ Phys.\ \textbf{67} (2019) 1900083,
\href{https://arxiv.org/abs/1906.00392}{arXiv:1906.00392}.

\bibitem{SuppFigures}
B.S.~Hartshorn,
\emph{Supplementary material: figures for the pyCICY-X package},
companion document (Pre-print 2026), \url{https://doi.org/10.5281/zenodo.21843383}.

\bibitem{AndersonGaoGrayLee2017}
L.~B.~Anderson, X.~Gao, J.~Gray and S.-J.~Lee,
\emph{Fibrations in CICY threefolds},
JHEP \textbf{10} (2017) 077,
\href{https://arxiv.org/abs/1708.07907}{arXiv:1708.07907}.

\bibitem{AndersonFibrations2018}
L.~B.~Anderson, J.~Gray and B.~Hammack,
\emph{Fibrations in non-simply connected Calabi--Yau quotients},
JHEP \textbf{08} (2018) 128,
\href{https://arxiv.org/abs/1805.05497}{arXiv:1805.05497}.

\bibitem{HKTY1994}
S.~Hosono, A.~Klemm, S.~Theisen and S.-T.~Yau,
\emph{Mirror symmetry, mirror map and applications to complete intersection
Calabi--Yau spaces},
Nucl.\ Phys.\ B \textbf{433} (1995) 501,
\href{https://arxiv.org/abs/hep-th/9406055}{arXiv:hep-th/9406055}.

\bibitem{Braun2010}
V.~Braun,
\emph{On free quotients of complete intersection Calabi--Yau manifolds},
JHEP \textbf{04} (2011) 005,
\href{https://arxiv.org/abs/1003.3235}{arXiv:1003.3235}.

\bibitem{ConstantinGrayLukas2017}
A.~Constantin, J.~Gray and A.~Lukas,
\emph{Hodge numbers for all CICY quotients},
JHEP \textbf{01} (2017) 001,
\href{https://arxiv.org/abs/1607.01830}{arXiv:1607.01830}.

\bibitem{CandelasConstantinMishra2016}
P.~Candelas, A.~Constantin and C.~Mishra,
\emph{Hodge numbers for CICYs with symmetries of order divisible by 4},
Fortsch.\ Phys.\ \textbf{64} (2016) 463,
\href{https://arxiv.org/abs/1511.01103}{arXiv:1511.01103}.

\bibitem{GrayWang2022}
J.~Gray and J.~Wang,
\emph{Free quotients of favorable Calabi--Yau manifolds},
JHEP \textbf{07} (2022) 116,
\href{https://arxiv.org/abs/2112.12683}{arXiv:2112.12683}.

\bibitem{LukasMishra2020}
A.~Lukas and C.~Mishra,
\emph{Discrete symmetries of complete intersection Calabi--Yau manifolds},
Commun.\ Math.\ Phys.\ \textbf{379} (2020) 847,
\href{https://arxiv.org/abs/1708.08943}{arXiv:1708.08943}.

\bibitem{TianYau1987}
G.~Tian and S.-T.~Yau,
\emph{Three dimensional algebraic manifolds with $c_1 = 0$ and $\chi = -6$},
in \emph{Mathematical Aspects of String Theory}, World Scientific (1987) 543.

\bibitem{CandelasIII}
P.~Candelas, C.~A.~L\"utken and R.~Schimmrigk,
\emph{Complete intersection Calabi--Yau manifolds II: three generation
manifolds},
Nucl.\ Phys.\ B \textbf{306} (1988) 113.

\bibitem{Brittenham2025}
M.~Brittenham and S.~Hermiller,
\emph{Unknotting number is not additive under connected sum},
\href{https://arxiv.org/abs/2506.24088}{arXiv:2506.24088}.

\bibitem{Wang2025}
Chao Wang, Yimu Zhang (2025)
\emph{A remark on the counterexample to the unknotting number conjecture}
\url{https://doi.org/10.48550/arXiv.2507.14265}


\bibitem{Brittenham2021}
M.~Brittenham and S.~Hermiller,
\emph{A counterexample to the Bernhard--Jablan Unknotting Conjecture},
Exp.\ Math.\ \textbf{30} (2021) 547.


\bibitem{AndersonGrayLukasPalti2011}
L.B.~Anderson, J.~Gray, A.~Lukas and E.~Palti,
\emph{Heterotic line bundle standard models},
Phys.\ Rev.\ D \textbf{84} (2011) 106005,
\href{https://arxiv.org/abs/1106.4804}{arXiv:1106.4804}.

\bibitem{Braun2011}
V.~Braun,
\emph{The 24-cell and Calabi--Yau threefolds with Hodge numbers $(1,1)$},
JHEP \textbf{05} (2012) 101,
\href{https://arxiv.org/abs/1102.4880}{arXiv:1102.4880}.

\bibitem{MarraProiettiSheng2024}
P.~Marra, V.~Proietti and X.~Sheng,
\emph{Hofstadter--Toda spectral duality and quantum groups},
J.\ Math.\ Phys.\ \textbf{65} (2024) 072102,
\href{https://arxiv.org/abs/2312.14242}{arXiv:2312.14242}.


\bibitem{Larfors2022}
M.~Larfors, A.~Lukas, F.~Ruehle and R.~Schneider,
\emph{Learning size and shape of Calabi--Yau spaces},
\href{https://arxiv.org/abs/2111.01436}{arXiv:2111.01436};
code at \url{https://github.com/pythoncymetric/cymetric}.

\bibitem{HendiLarforsWalden2024}
A.~Hendi, M.~Larfors and M.~Walden,
\emph{Learning Group Invariant Calabi--Yau Metrics by Fundamental Domain
Projections},
\href{https://arxiv.org/abs/2407.06914}{arXiv:2407.06914}.

\bibitem{Butbaia2024}
G.~Butbaia, D.~Mayorga Pe\~na, J.~Tan, P.~Berglund, T.~H\"ubsch,
V.~Jejjala and C.~Mishra,
\emph{Physical Yukawa couplings in heterotic string compactifications},
\href{https://arxiv.org/abs/2401.15078}{arXiv:2401.15078};
code at \url{https://github.com/Justin-Tan/cymyc}.

\bibitem{Anderson2023branes}
L.~B.~Anderson, C.~R.~Brodie and J.~Gray,
\emph{Branes and bundles through conifold transitions and dualities in heterotic string theory},
Phys.\ Rev.\ D \textbf{108} (2023) 106018,
\href{https://arxiv.org/abs/2211.05804}{arXiv:2211.05804}.




\bibitem{Green1987}
P.~Green and T.~H\"ubsch,
\emph{Calabi--Yau manifolds as complete intersections in products of complex projective spaces},
Commun.\ Math.\ Phys.\ \textbf{109} (1987) 99.

\bibitem{Green1989}
P.~S.~Green, T.~H\"ubsch and C.~A.~L\"utken,
\emph{All the Hodge Numbers for All Calabi--Yau Complete Intersections},
Class.\ Quant.\ Grav.\ \textbf{6} (1989) 105.


\bibitem{Aharonov2023}
Y.~Aharonov, D.~Collins and S.~Popescu,
\emph{Angular momentum flows without anything carrying it},
\href{https://arxiv.org/abs/2310.07568}{arXiv:2310.07568}.

\bibitem{Waegell2024}
M.~Waegell, J.~Tollaksen and Y.~Aharonov,
\emph{Separating a particle's mass from its momentum},
Quantum \textbf{8} (2024) 1536,
\href{https://arxiv.org/abs/2401.10408}{arXiv:2401.10408}.



\bibitem{GrayHauptLukas2013}
J.~Gray, A.~S.~Haupt and A.~Lukas,
\emph{All complete intersection Calabi--Yau four-folds},
JHEP \textbf{07} (2013) 070,
\href{https://arxiv.org/abs/1303.1832}{arXiv:1303.1832}.


\bibitem{Anderson2011}
L.~B.~Anderson, J.~Gray, A.~Lukas and E.~Palti,
\emph{Two hundred heterotic standard models on smooth Calabi--Yau threefolds},
Phys.\ Rev.\ D \textbf{84} (2011) 106005,
\href{https://arxiv.org/abs/1106.4804}{arXiv:1106.4804}.

\bibitem{Anderson2012}
L.~B.~Anderson, J.~Gray, A.~Lukas and E.~Palti,
\emph{Heterotic line bundle standard models},
JHEP \textbf{06} (2012) 113,
\href{https://arxiv.org/abs/1202.1757}{arXiv:1202.1757}.





\end{thebibliography}

\end{document}
